\documentclass[12pt]{article}
\usepackage{amsmath, amssymb, amsthm}
\usepackage{geometry}
\usepackage{hyperref}
\usepackage{listings}
\usepackage{courier}
\usepackage{setspace}
\usepackage{titlesec}

\geometry{margin=1in}
\setstretch{1.15}

\titleformat{\section}{\large\bfseries}{\thesection}{0.5em}{}

\lstset{
    basicstyle=\footnotesize\ttfamily,
    breaklines=true,
    frame=single,
    columns=fullflexible
}

\title{Exploratory Thoughts: Unified Recommender System Architecture\\
\large Transformer Sequence Modeling + Feature Interaction\\
\normalsize \href{https://algo.qq.com/\#intro}{algo.qq.com}}
\author{David \and Xia}
\date{April 2026 \\ \small Draft v0.02}

\begin{document}
\maketitle

\begin{center}
    \small\textit{Prepared in the context of the QQ Algorithm Challenge ---
    \href{https://algo.qq.com/\#intro}{algo.qq.com}}
\end{center}

\begin{abstract}
These are first-thoughts towards a unified recommender system architecture
that merges transformer-based sequence modelling with structured feature
interaction, targeting the QQ Algorithm Challenge. This version explores a
TensorFlow/Keras implementation. The intent is exploratory: to sketch the
mathematical and implementation foundations before committing to a full
proposal.
\end{abstract}

\noindent\textbf{Keywords:} recommender systems, transformers, feature interaction, TensorFlow, Keras, deep learning

\bigskip

\section{Overview}
This document presents a compact technical model describing how a unified
recommender-system architecture can merge:
\begin{itemize}
    \item Transformer-style sequence modeling,
    \item Feature-interaction modeling (e.g.\ DeepFM, DCN, DIN),
    \item A unified TensorFlow/Keras implementation,
    \item A training pipeline using the Keras \texttt{fit} API.
\end{itemize}

The goal is to show how modern recommender systems integrate temporal
dependencies and structured feature interactions into a single scalable model,
leveraging TensorFlow's production-ready ecosystem.

\section{1.\ Minimal Unified Architecture (TensorFlow/Keras)}
Below is a minimal Keras model that combines:
\begin{itemize}
    \item \texttt{MultiHeadAttention} for sequential user behaviour,
    \item static feature embeddings (user, item, context),
    \item an MLP for feature interactions.
\end{itemize}

\subsection*{Model Code}
\begin{lstlisting}[language=Python]
import tensorflow as tf
from tensorflow import keras
from tensorflow.keras import layers

class UnifiedRecModel(keras.Model):
    def __init__(
        self,
        num_items,
        num_users,
        num_context_feats,
        d_model=64,
        n_heads=4,
        d_ff=128,
        mlp_hidden=128,
    ):
        super().__init__()

        self.item_emb = layers.Embedding(num_items, d_model)
        self.user_emb = layers.Embedding(num_users, d_model)
        self.ctx_emb  = layers.Embedding(num_context_feats, d_model)

        self.attention = layers.MultiHeadAttention(
            num_heads=n_heads,
            key_dim=d_model // n_heads,
        )
        self.ffn = keras.Sequential([
            layers.Dense(d_ff, activation="relu"),
            layers.Dense(d_model),
        ])
        self.norm1 = layers.LayerNormalization()
        self.norm2 = layers.LayerNormalization()

        self.mlp = keras.Sequential([
            layers.Dense(mlp_hidden, activation="relu"),
            layers.Dense(1),
        ])

    def call(self, inputs, training=False):
        seq_items, user_ids, target_items, ctx_feats, seq_mask = inputs

        seq_emb  = self.item_emb(seq_items)          # (B, T, d)
        user_vec = self.user_emb(user_ids)            # (B, d)
        item_vec = self.item_emb(target_items)        # (B, d)
        ctx_vec  = tf.reduce_mean(
            self.ctx_emb(ctx_feats), axis=1
        )                                             # (B, d)

        # self-attention with padding mask
        attn_mask = seq_mask[:, tf.newaxis, tf.newaxis, :]
        attn_out  = self.attention(
            seq_emb, seq_emb, attention_mask=attn_mask
        )
        seq_enc = self.norm1(seq_emb + attn_out)
        seq_enc = self.norm2(seq_enc + self.ffn(seq_enc))

        # masked mean-pool
        mask_f   = tf.cast(seq_mask, tf.float32)
        lengths  = tf.reduce_sum(mask_f, axis=1, keepdims=True)
        seq_pool = tf.reduce_sum(
            seq_enc * mask_f[:, :, tf.newaxis], axis=1
        ) / tf.maximum(lengths, 1.0)

        x = tf.concat([seq_pool, user_vec, item_vec, ctx_vec], axis=-1)
        return tf.squeeze(self.mlp(x), axis=-1)
\end{lstlisting}

\section{2.\ Merging Attention and Feature-Interaction Layers}
A unified architecture combines:
\begin{itemize}
    \item Transformer attention for temporal modelling,
    \item Feature-interaction layers (e.g.\ FM, DeepFM) for structured fields.
\end{itemize}

\subsection*{Attention}
\[
\mathrm{Attn}(Q,K,V)
= \mathrm{softmax}\!\left(\frac{QK^\top}{\sqrt{d_k}}\right)V.
\]

\subsection*{Factorization Machine Interaction}
\[
y_{\mathrm{FM}}
= \frac{1}{2}\sum_{i<j}
\left(
\langle v_i, v_j\rangle
\right).
\]

\subsection*{FM Layer Code}
\begin{lstlisting}[language=Python]
class FMInteraction(keras.layers.Layer):
    def call(self, field_embs):
        # field_embs: (B, num_fields, d)
        sum_emb = tf.reduce_sum(field_embs, axis=1)
        sum_sq  = sum_emb * sum_emb
        sq_sum  = tf.reduce_sum(field_embs * field_embs, axis=1)
        fm = 0.5 * tf.reduce_sum(sum_sq - sq_sum, axis=1)
        return fm
\end{lstlisting}

\section{3.\ Training Pipeline}
TensorFlow's \texttt{model.compile}/\texttt{model.fit} API provides a concise
training pipeline:
\begin{enumerate}
    \item Data preparation (sequence padding, masks, feature IDs) as \texttt{tf.data.Dataset}.
    \item Compile with Adam optimiser and binary cross-entropy loss.
    \item Call \texttt{model.fit} with built-in AUC metric tracking.
    \item Export via TF-Serving or TFLite for deployment.
\end{enumerate}

\subsection*{Training Loop}
\begin{lstlisting}[language=Python]
model = UnifiedRecModel(
    num_items, num_users, num_context_feats
)

model.compile(
    optimizer=tf.keras.optimizers.Adam(learning_rate=1e-3),
    loss=tf.keras.losses.BinaryCrossentropy(from_logits=True),
    metrics=[tf.keras.metrics.AUC(name="auc")],
)

model.fit(
    train_dataset,
    validation_data=val_dataset,
    epochs=num_epochs,
    callbacks=[
        tf.keras.callbacks.EarlyStopping(
            monitor="val_auc",
            patience=3,
            restore_best_weights=True,
        )
    ],
)
\end{lstlisting}

\section{4.\ TensorFlow Ecosystem and Deployment}
\begin{itemize}
    \item \textbf{TF-Serving}: production-grade REST/gRPC model serving with versioning.
    \item \textbf{TFLite}: on-device inference for mobile and edge scenarios.
    \item \textbf{TF Extended (TFX)}: end-to-end ML pipelines with data validation and model monitoring.
    \item \textbf{SavedModel format}: language-agnostic serialisation for cross-platform deployment.
    \item \textbf{Trend}: TensorFlow remains the dominant choice in large-scale industry CTR/CVR systems.
\end{itemize}

\end{document}
